\chapter{Discussion, Limitations and Threats to Validity}
\label{sec:discussion}

% Source notes: Project/docs/thesis_notes.md, checkpoint_2026-10-10.md (takeaways and caveats).

This chapter interprets the results along the five research questions, then states what the experiments cannot show
and which choices could have biased them.

\section{When a specialist decision model pays off}
\label{sec:disc-synthesis}

\paragraph{Specialisation (RQ1).} Training on the decision is what makes GroundOne accurate. The untrained models are
near chance, and the same models after training reach 0.769 and 0.832 three-way balanced accuracy on the RAGTruth test
set, beating the specialised encoders in AUROC and Rizzo-Flow, a general-purpose decision model eight times larger, in three-way accuracy.
The advantage is concentrated where the training data are: on LLM-AggreFact the 3B is on par with the best small encoder
rather than above it, and below the closed decision model Jev and a 27B open model as measured independently; the 0.5B
remains below the encoder, and on VitaminC the trained 3B is no better than the untrained one. The comparison with d1-3B, a general-purpose decision model of the same size, makes the trade-off explicit. d1-3B, already
trained to decide but on a broad mixture of tasks, is far behind GroundOne on RAGTruth, especially on long answers, on par
with the 3B on LLM-AggreFact and far ahead on VitaminC. A specialist is therefore a strong choice for the domain it was
trained on and a weak general fact checker: specialisation buys accuracy and robustness on long answers in its domain at
the cost of transfer, while model size, or a broad training mixture, is what improves transfer outside it.

\paragraph{Readout versus generation (RQ2).} The decision-model mechanism contributes in two separable ways. The first is
cost: one pass over the answer replaces the generation of $K$ verdicts, so the time saved grows with the length of the
answer, reaching a factor of about two against a terse judge and six against a verbose one for answers with fourteen or
more sentences. The second, which this thesis did not expect, is accuracy. With identical weights and cost, deciding
among the three admissible letters is more accurate than taking the token the model would generate, by 0.25 for the
0.5B and 0.07 for the 3B. A fine-tuned decision model is better understood as a three-way classifier placed on top of a
language model than as a generator that stops after one token, and this distinction matters most for small models,
which are the ones a laptop can afford. The comparison with generative judges, finally, also measures training: the
judges are untrained, so the large differences in accuracy between them and GroundOne should be read as the combined
effect of specialisation and readout.

\paragraph{Calibration (RQ3).} A temperature and a bias fitted on development data make the probabilities match the
observed frequencies on the test set of the same domain, and the calibrated probabilities make selective prediction
effective. Two lessons go beyond GroundOne. First, the expected calibration error alone cannot rank systems: after the
same calibration, judges whose probabilities barely separate the classes reach the same ECE as GroundOne, and only the
AUROC and the risk--coverage curve reveal the difference. Second, decisions and probabilities must be derived
consistently: the most probable calibrated class collapses towards the majority class, while a threshold fitted on
development data recovers the balanced decision and keeps the label consistent with the probability shown to the user.
Calibration is valid only for the distribution it was fitted on, as the results on VitaminC show.

\paragraph{Cost on the device (RQ4).} The verifier fits comfortably in the 16\,GB laptop: training peaks at 2.6 and
5.9\,GB, and verifying a long answer takes 0.41 seconds with the 0.5B and 2.57 with the 3B, at an energy of 5.5 and 33
joules. The 0.5B meets an interactive budget; the 3B is slower than an interactive target for a typical answer and
trades time for accuracy. The energy measurements qualify the speed results: a single long prefill draws more power than
token-by-token decoding, so the energy saved is smaller than the time saved, and speed should not be read as an equal
saving of battery. The verifier also coexists with a generator in the same process: together with a 4-bit 3B
generator, the peak memory is 6.1\,GB with the 3B and 3.3\,GB with the 0.5B. The whole evaluation runs without any
network access, with results identical bit for bit, and 8-bit weights cost no accuracy compared with bf16 ones
(Section~\ref{sec:results-locality}). The retriever remains to be measured, which will be done in the demonstration
application.

The economic cost should also be read with care (Table~\ref{tab:disc-cost}). Verifying 10{,}000 answers with GroundOne
takes an amount of energy that, at the European average household electricity price~\cite{eurostat2026electricity},
is worth less than one euro cent with the 0.5B and about three cents with the 3B. Entrusting the same verification to a
hosted judge, at the providers' list prices~\cite{anthropic2026pricing,openai2026pricing,google2026geminipricing},
costs instead from about one dollar with the cheapest models to more than two hundred with the top-tier ones. The ratio
is thus two to four orders of magnitude, but in absolute terms cheap judges remain inexpensive, and the comparison
includes neither the purchase of the hardware nor the quality of hosted judges, which this thesis did not measure. The
decisive argument for local verification therefore remains independence: the documents never leave the machine and the
service does not depend on an external provider.

\begin{table}[htbp]
\centering
\caption{Indicative cost of verifying 10{,}000 RAG answers. Hosted judges: standard-tier list prices, accessed on 19
October 2026, with about 900 input tokens per answer (median T1 document and answer plus instructions) and the output
tokens measured for the 3B judges of this thesis (36 for the terse judge, 248 for the verbose one). GroundOne: measured
net chip energy multiplied by the EU average household electricity price (0.29~€/kWh, second half of 2025). Hardware
cost and caching or batch discounts are excluded.}
\label{tab:disc-cost}
\small
\begin{tabular}{lccc}
\toprule
System & \makecell{Input / output price\\(USD per Mtoken)} & \makecell{Terse\\judge} & \makecell{Verbose\\judge} \\
\midrule
\makecell[l]{Claude Haiku 5.5,\\gpt-6-luna} & 0.10 / 0.50 & 1.1 USD & 2.2 USD \\
Gemini 3.5 Flash-Lite & 0.30 / 2.50 & 3.6 USD & 8.9 USD \\
Gemini 3.8 Flash & 0.75 / 3.75 & 8.2 USD & 16 USD \\
\makecell[l]{Claude Sonnet 5.5, gpt-6.1-sol,\\Gemini 4 Argon} & 2.00 / 10.00 & 22 USD & 43 USD \\
Claude Opus 5.5 & 4.00 / 20.00 & 43 USD & 86 USD \\
gpt-6-astra & 10.00 / 50.00 & 109 USD & 215 USD \\
\midrule
GroundOne 0.5B (local) & \multicolumn{3}{c}{0.015 kWh, about 0.004 €} \\
GroundOne 3B (local) & \multicolumn{3}{c}{0.092 kWh, about 0.03 €} \\
\bottomrule
\end{tabular}
\end{table}

\paragraph{Beyond a binary score (RQ5).} The three-way output separates contradicted from not-mentioned sentences well
in data-to-text, where RAGTruth concentrates its contradictions, and much less in question answering and summarisation,
where contradictions are rare in the training data. The 3B improves this considerably, which suggests that more varied
contradictions in training, rather than a different architecture, are what the small model lacks. Because the strict and
lenient policies and the threshold are derived from the same three probabilities, a user can change the notion of error
without retraining, which no binary verifier allows.

\section{Limitations}
\label{sec:disc-limitations}

\paragraph{One platform and one family of models.} Every result comes from one laptop, one runtime and one model
family. Latency and energy depend on the hardware and on the implementation, and the ratios between methods, not the
absolute numbers, are what should transfer. No latency comparison is made with the encoders or the general-purpose System
One, because they ran in different runtimes and under conditions that would make the comparison unfair.

\paragraph{Training stability.} Two completed runs of the 0.5B agree on all aggregate metrics, but a third run collapsed
early and was stopped. The recipe has no gradient clipping and uses Adam without bias correction, both known causes of
unstable fine-tuning~\cite{pascanu2013difficulty,mosbach2021stability}, so either may explain the collapse; neither was
tested. The 3B was trained once, on half of the training answers, and its seed variance is unknown.

\paragraph{Training data.} RAGTruth is the only training source. Its contradictions are concentrated in data-to-text,
and outside that task they are learned only in part. The models also rely partly on lexical overlap: paraphrased support
looks unsupported, and a small change in a sentence that otherwise copies the source often goes unnoticed. This blind
spot shrinks with model size but remains the main limitation of both models.

\paragraph{Scope of the decision.} Verification is sentence by sentence and only against the retrieved document. A
sentence that is true in the world but absent from the document counts as not mentioned, which is the right behaviour for
groundedness but not for fact checking. Sentences without checkable content are excluded by a rule, and the segmentation
itself can split or merge claims differently from a human reader. The verdict on a sentence depends on the document and on
the preceding sentences only.

\paragraph{Contexts made of several documents.} When the answer is written from several retrieved passages, GroundOne
reads them concatenated as a single document, as in RAGTruth question answering. Three limits remain. The first concerns
length: in training, document and answer never exceed 2{,}048 tokens, so a longer context falls outside the distribution
on which the model was measured and must be split into chunks, keeping for each sentence the highest support probability,
as on LLM-AggreFact; a fact that requires combining two different chunks may then be missed. The second concerns
conflicting sources: when one passage supports a sentence and another contradicts it, the expected verdict is not
uniquely defined, and RAGTruth contains no such cases. The third concerns attribution: GroundOne states whether a sentence
is supported, but not by which passage.

\paragraph{Statistical scope.} The bootstrap intervals describe the variability of the test data, not that of training,
and about fifty comparisons were tested without correction for multiple testing. Conclusions therefore rest on effects far
from zero, and differences on LLM-AggreFact and VitaminC smaller than the spread between seeds are not claimed. Several
baselines were evaluated on subsamples, which widens their intervals.

\section{Threats to validity}
\label{sec:disc-threats}

\paragraph{Labels.} The distinction between contradicted and not mentioned rests on the RAGTruth annotators, whose
agreement on the type of hallucination is not published, and a recent re-annotation of the RAGTruth test split found
considerably more hallucinated answers than the original labels~\cite{yu2026retromorphic}. The test set may therefore
under-count the minority classes. For this reason the three-way results on RAGTruth are always read together with those
on VitaminC, whose labels have measured agreement. LLM-AggreFact unifies the labels of eleven datasets of different
quality, and some of its columns are known to be noisy.

\paragraph{In-domain advantage.} GroundOne is trained on RAGTruth, so its results on the RAGTruth test set and on the
RAGTruth column of LLM-AggreFact favour it with respect to systems trained elsewhere; conversely, VitaminC lies in the
training lineage of MiniCheck-RoBERTa. Each comparison is reported with the relevant caveat, and the claims rest on the
pattern across the three test sets rather than on a single one.

\paragraph{Baseline configuration.} Generative judges are sensitive to their prompts. Their prompts were fixed before
looking at test results, and the prompt of the equal-output judge was chosen on development data for format validity
only; parsing was lenient in their favour; decoding was not constrained by a grammar, which mostly affects the format
failures of the smallest model. A judge fine-tuned on the same data, or a larger judge, was not evaluated, so the
comparison measures what a judge of the same size gives without training. Rizzo-Flow was run with
4-bit weights because the 8-bit model did not fit beside the operating system, which may understate its accuracy. The
d1-3B was run zero-shot with the questions used for all decision-model baselines, which may not
be the phrasing it works best with, and on a subsample of LLM-AggreFact.

\paragraph{Contamination.} The base models were pretrained on large web corpora that may include parts of the public
test sets. Because the same base models are used untrained as baselines and reach near-chance accuracy on the
three-way task, contamination cannot explain the gains of training, but it cannot be ruled out entirely for the absolute
numbers.

\paragraph{One machine.} All measurements of time and energy come from one laptop, an Apple M4 with 16\,GB, and one
runtime, MLX. The cost model of Chapter~\ref{sec:local} predicts that the advantage of a decision model holds on every class
of local hardware, because it rests on the difference between prefill and decoding, which vendors and independent studies
observe everywhere; the lower bounds on decoding computed there for other machines follow from published specifications.
The actual speed-ups on a discrete graphics card, a DGX Spark or a phone were not measured: their size depends on the
balance between computation and bandwidth of each machine, and the direction of the advantage, not its magnitude, is
what the cost model supports.

